% 本実験3（時間重み付け比較）の詳細（equal以外を含む）

本文（本実験3）では、議論を明確にするため特徴量配分を \texttt{equal} に固定した結果のみを示した。
本付録では、\textbf{他の特徴量配分を含む}時間重み付け比較の結果をまとめる。

\begin{table}[htbp]
\centering
\caption{時間重み付け比較（全特徴量配分平均, 複合スコア比較）}
\scriptsize
\begin{tabular}{l r r r r l}
\hline
代表パターン & 複合スコア & 平均Top-1 & 平均Top-3 & 平均Top-5 & 備考 \\
\hline
直近重視の指数減衰 & 0.6199 & 0.5352 & 0.6880 & 0.7296 & 全重み付け平均で最良 \\
直近30日線形減衰 & 0.6131 & 0.5278 & 0.6824 & 0.7222 & 線形に直近を重視 \\
時間減衰なし & 0.5734 & 0.4917 & 0.6278 & 0.6963 & ベースライン \\
\hline
\end{tabular}
\end{table}

\begin{table}[htbp]
\centering
\caption{時間重み付け上位パターン（学習20件, 全特徴量配分）}
\scriptsize
\begin{tabularx}{\linewidth}{X X r r r}
\hline
特徴量配分 & 時間重み & 平均Top-1 & 平均Top-3 & 平均Top-5 \\
\hline
業務内容重視（\texttt{content\_focused}） & 直近重視の指数減衰 & 0.5111 & 0.6889 & 0.7500 \\
業務内容重視（\texttt{content\_focused}） & セッション順の線形減衰（短期） & 0.5111 & 0.7000 & 0.7500 \\
均等配分（\texttt{equal}） & セッション順の線形減衰（中期） & 0.5111 & 0.6333 & 0.7444 \\
\hline
\end{tabularx}
\end{table}

\begin{table}[htbp]
\centering
\caption{時間重み付け上位パターン（学習20件, 全特徴量配分）}
\scriptsize
\begin{tabularx}{\linewidth}{X X r r r}
\hline
特徴量配分 & 時間重み & 平均Top-1 & 平均Top-3 & 平均Top-5 \\
\hline
業務内容重視（\texttt{content\_focused}） & セッション順線形10件（\texttt{time\_session\_linear\_10}） & 0.5111 & 0.7000 & 0.7500 \\
業務内容重視（\texttt{content\_focused}） & 直近7日指数減衰（\texttt{time\_recent\_exp\_7d}） & 0.5111 & 0.6889 & 0.7500 \\
均等配分（\texttt{equal}） & セッション順線形20件（\texttt{time\_session\_linear\_20}） & 0.5111 & 0.6333 & 0.7444 \\
業務内容重視（\texttt{content\_focused}） & Softmax 7日（\texttt{time\_softmax\_7d}） & 0.4778 & 0.6889 & 0.7444 \\
業務内容重視（\texttt{content\_focused}） & 直近30日線形減衰（\texttt{time\_recent\_linear\_30d}） & 0.5278 & 0.7000 & 0.7389 \\
時給重視（\texttt{salary\_focused}） & セッション順線形20件（\texttt{time\_session\_linear\_20}） & 0.5167 & 0.6389 & 0.7389 \\
勤務時間重視（\texttt{work\_time\_focused}） & セッション順線形20件（\texttt{time\_session\_linear\_20}） & 0.5222 & 0.6222 & 0.7389 \\
均等配分（\texttt{equal}） & セッション順線形10件（\texttt{time\_session\_linear\_10}） & 0.5389 & 0.6667 & 0.7389 \\
均等配分（\texttt{equal}） & 直近7日指数減衰（\texttt{time\_recent\_exp\_7d}） & 0.5556 & 0.6944 & 0.7389 \\
均等配分（\texttt{equal}） & 直近90日線形減衰（\texttt{time\_recent\_linear\_90d}） & 0.5278 & 0.6611 & 0.7333 \\
均等配分（\texttt{equal}） & セッション順指数（半減期5）（\texttt{time\_session\_exp\_5}） & 0.5222 & 0.6611 & 0.7333 \\
業務内容重視（\texttt{content\_focused}） & 直近5件ステップ（\texttt{time\_step\_k\_5}） & 0.5333 & 0.7056 & 0.7333 \\
業務内容重視（\texttt{content\_focused}） & セッション順ステップ5件（\texttt{time\_session\_step\_k\_5}） & 0.5333 & 0.7056 & 0.7333 \\
介護項目重視（\texttt{care\_focused}） & 直近7日指数減衰（\texttt{time\_recent\_exp\_7d}） & 0.5500 & 0.6944 & 0.7333 \\
勤務時間重視（\texttt{work\_time\_focused}） & セッション順線形10件（\texttt{time\_session\_linear\_10}） & 0.5389 & 0.6611 & 0.7278 \\
時給重視（\texttt{salary\_focused}） & 直近30日指数減衰（\texttt{time\_recent\_exp\_30d}） & 0.5278 & 0.6556 & 0.7278 \\
時給重視（\texttt{salary\_focused}） & セッション順指数（半減期10）（\texttt{time\_session\_exp\_10}） & 0.5056 & 0.6444 & 0.7278 \\
業務内容重視（\texttt{content\_focused}） & 直近30日指数減衰（\texttt{time\_recent\_exp\_30d}） & 0.5278 & 0.6667 & 0.7278 \\
業務内容重視（\texttt{content\_focused}） & Softmax 30日（\texttt{time\_softmax\_30d}） & 0.5222 & 0.6667 & 0.7278 \\
均等配分（\texttt{equal}） & 直近30日指数減衰（\texttt{time\_recent\_exp\_30d}） & 0.5389 & 0.6556 & 0.7278 \\
均等配分（\texttt{equal}） & \texttt{time\_bi\_level\_10} & 0.5000 & 0.6333 & 0.7278 \\
均等配分（\texttt{equal}） & \texttt{time\_session\_bi\_level\_10} & 0.5000 & 0.6333 & 0.7278 \\
勤務時間重視（\texttt{work\_time\_focused}） & セッション順指数（半減期10）（\texttt{time\_session\_exp\_10}） & 0.5111 & 0.6278 & 0.7278 \\
介護項目重視（\texttt{care\_focused}） & 直近30日線形減衰（\texttt{time\_recent\_linear\_30d}） & 0.5389 & 0.6889 & 0.7278 \\
均等配分（\texttt{equal}） & 直近30日線形減衰（\texttt{time\_recent\_linear\_30d}） & 0.5333 & 0.6722 & 0.7278 \\
\hline
\multicolumn{5}{l}{（参考）直近90日ウィンドウ（\texttt{time\_window\_uniform\_90d}）の結果}\\
業務内容重視（\texttt{content\_focused}） & 直近90日ウィンドウ（\texttt{time\_window\_uniform\_90d}） & 0.4944 & 0.6500 & 0.7167 \\
均等配分（\texttt{equal}） & 直近90日ウィンドウ（\texttt{time\_window\_uniform\_90d}） & 0.4944 & 0.6500 & 0.7056 \\
時給重視（\texttt{salary\_focused}） & 直近90日ウィンドウ（\texttt{time\_window\_uniform\_90d}） & 0.4889 & 0.6500 & 0.7056 \\
求人タイトル重視（\texttt{job\_title\_focused}） & 直近90日ウィンドウ（\texttt{time\_window\_uniform\_90d}） & 0.4278 & 0.6444 & 0.6889 \\
介護項目重視（\texttt{care\_focused}） & 直近90日ウィンドウ（\texttt{time\_window\_uniform\_90d}） & 0.5111 & 0.6556 & 0.6778 \\
勤務時間重視（\texttt{work\_time\_focused}） & 直近90日ウィンドウ（\texttt{time\_window\_uniform\_90d}） & 0.5056 & 0.6556 & 0.6944 \\
\hline
\end{tabularx}

\end{table}

\begin{table}[htbp]
\centering
\caption{時間重み付け上位パターン（学習30件, 全特徴量配分）}
\scriptsize
\begin{tabularx}{\linewidth}{X X r r r}
\hline
特徴量配分 & 時間重み & 平均Top-1 & 平均Top-3 & 平均Top-5 \\
\hline
介護項目重視（\texttt{care\_focused}） & 直近30日ウィンドウ & 0.4182 & 0.5636 & 0.6909 \\
勤務時間重視（\texttt{work\_time\_focused}） & 直近5件ステップ & 0.4455 & 0.5727 & 0.6818 \\
勤務時間重視（\texttt{work\_time\_focused}） & セッション順ステップ5件 & 0.4455 & 0.5727 & 0.6818 \\
介護項目重視（\texttt{care\_focused}） & 直近5件ステップ & 0.4364 & 0.5545 & 0.6818 \\
介護項目重視（\texttt{care\_focused}） & セッション順ステップ5件 & 0.4364 & 0.5545 & 0.6818 \\
時給重視（\texttt{salary\_focused}） & 直近30日ウィンドウ & 0.3455 & 0.5364 & 0.6818 \\
業務内容重視（\texttt{content\_focused}） & 直近5件ステップ & 0.4455 & 0.6182 & 0.6727 \\
業務内容重視（\texttt{content\_focused}） & セッション順ステップ5件 & 0.4455 & 0.6182 & 0.6727 \\
勤務時間重視（\texttt{work\_time\_focused}） & 直近30日ウィンドウ & 0.4545 & 0.5727 & 0.6727 \\
求人タイトル重視（\texttt{job\_title\_focused}） & 直近5件ステップ & 0.4273 & 0.5636 & 0.6636 \\
求人タイトル重視（\texttt{job\_title\_focused}） & セッション順ステップ5件 & 0.4273 & 0.5636 & 0.6636 \\
均等配分（\texttt{equal}） & 直近30日ウィンドウ & 0.3636 & 0.5364 & 0.6636 \\
時給重視（\texttt{salary\_focused}） & 直近5件ステップ & 0.4091 & 0.5636 & 0.6545 \\
時給重視（\texttt{salary\_focused}） & セッション順ステップ5件 & 0.4091 & 0.5636 & 0.6545 \\
均等配分（\texttt{equal}） & 直近5件ステップ & 0.4273 & 0.5545 & 0.6545 \\
均等配分（\texttt{equal}） & セッション順ステップ5件 & 0.4273 & 0.5545 & 0.6545 \\
介護項目重視（\texttt{care\_focused}） & 直近30日線形減衰 & 0.4182 & 0.5364 & 0.6545 \\
業務内容重視（\texttt{content\_focused}） & S字型減衰（直近重視） & 0.4182 & 0.5273 & 0.6545 \\
均等配分（\texttt{equal}） & 直近30日線形減衰 & 0.4273 & 0.5273 & 0.6455 \\
求人タイトル重視（\texttt{job\_title\_focused}） & 直近30日線形減衰 & 0.4091 & 0.5273 & 0.6455 \\
時給重視（\texttt{salary\_focused}） & 直近30日線形減衰 & 0.3909 & 0.5273 & 0.6455 \\
業務内容重視（\texttt{content\_focused}） & 直近30日ウィンドウ & 0.4455 & 0.5455 & 0.6364 \\
業務内容重視（\texttt{content\_focused}） & 直近30日線形減衰 & 0.4000 & 0.5455 & 0.6364 \\
求人タイトル重視（\texttt{job\_title\_focused}） & Softmax 7日 & 0.4091 & 0.5364 & 0.6364 \\
求人タイトル重視（\texttt{job\_title\_focused}） & 直近7日指数減衰 & 0.4182 & 0.5273 & 0.6364 \\
\hline
\multicolumn{5}{l}{（参考）直近90日ウィンドウの結果}\\
業務内容重視（\texttt{content\_focused}） & 直近90日ウィンドウ & 0.3818 & 0.4727 & 0.5455 \\
均等配分（\texttt{equal}） & 直近90日ウィンドウ & 0.3545 & 0.4636 & 0.5273 \\
時給重視（\texttt{salary\_focused}） & 直近90日ウィンドウ & 0.3364 & 0.4636 & 0.5364 \\
求人タイトル重視（\texttt{job\_title\_focused}） & 直近90日ウィンドウ & 0.3818 & 0.4727 & 0.5364 \\
介護項目重視（\texttt{care\_focused}） & 直近90日ウィンドウ & 0.3909 & 0.4727 & 0.5455 \\
勤務時間重視（\texttt{work\_time\_focused}） & 直近90日ウィンドウ & 0.3727 & 0.4727 & 0.5273 \\
\hline
\end{tabularx}

\end{table}
